\chapter{Zeitdiskrete Systeme im z-Bereich}\label{sec:zeitdiskrete-systeme-im-z-bereich}

In Kapitel \hyperref[z-transformation-von-signalen]{5} wird die z-Transformation von Signalen vorgestellt und an Beispielen vertieft. Die z-Transformation kann zur analytischen Lösung von Differenzengleichungen mit konstanten Koeffizienten eingesetzt werden. Es wird sich zeigen, dass die analytische Lösung von Differenzengleichungen mit Hilfe der z-Transformation einfacher und übersichtlicher ist als die Vier-Schritt-Methode im Zeitbereich.

Die Lösung von linearen Differenzengleichungen führt zu dem Begriff der Übertragungsfunktion für zeitdiskrete Systeme. Sie weist Parallelen zur Übertragungsfunktion zeitkontinuierlicher Systeme auf. Auch die Interpretation der Übertragungsfunktion ist ähnlich. So lassen sich zum Beispiel Stabilitätsaussagen und Aussagen zur Schwingungsneigung sowie zur Sprungfähigkeit an der Pollage der Übertragungsfunktion ablesen. Wesentlich dabei ist, den in Kapitel 5 beschriebenen Zusammenhang zwischen der s- und z-Ebene zu kennen.

Die Interpretation der unterschiedlichen Systeme führt zur Einteilung in Systeme mit endlicher und unendlicher Impulsantwort. Aufgrund ihrer Impulsantworten werden sie als Finite-Impulse-Response-Systeme (FIR-Systeme) beziehungsweise Infinite-Impulse-Response-Systeme (IIR-Systeme) bezeichnet. Bei dem Entwurf zeitdiskreter Systeme wird sich zeigen, dass FIR- und IIR-Systeme grundsätzlich unterschiedliche Entwurfsverfahren erfordern.

In der Praxis wird die Simulation und Interpretation zeitdiskreter Systeme mit Programmen wie MATLAB und Simulink durchgeführt. Die dazu erforderlichen Befehle und Methoden werden kurz vorgestellt und an Beispielen verdeutlicht.

\section{Lösung von Differenzengleichungen mit der z-Transformation}\label{sec:luxf6sung-von-differenzengleichungen-mit-der-z-transformation}

Differenzengleichungen beschreiben das Verhalten zeitdiskreter Systeme ähnlich, wie Differentialgleichungen das Verhalten zeitkontinuierlicher Systeme beschreiben. Differenzengleichungen können im Zeitbereich gelöst werden. Das Verfahren dazu wird in Kapitel \hyperref[_Toc241453879]{0} vorgestellt. Es besteht aus mehreren Schritten und ist insbesondere bei Berücksichtigung von Anfangsbedingungen aufwendig. Eine effizientere Möglichkeit zur Lösung von Differenzengleichungen wird mit der z-Transformation erreicht. Das Vorgehen wird in Bild \ref{fig:6.1} skizziert.

\begin{figure}
  \centering
  \includesvg{./media/6_1_verfahren_zur_lsung_von_differen.svg}
  \caption{Verfahren zur Lösung von Differenzengleichungen mit der z-Transformation}
  \label{fig:6.1}
\end{figure}

Die Differenzengleichung wird mit Hilfe der z-Transformation in den z-Bereich transformiert. Aufgrund der Rechenregeln der z-Transformation ergibt sich eine algebraische Gleichung, die im\\
z-Bereich gelöst wird. Mit Hilfe der Rücktransformation über Partialbruchzerlegung wird die Lösung im Zeitbereich bestimmt.

\subsection{Lösung linearer Differenzengleichungen ohne Anfangsbedingungen}\label{sec:luxf6sung-linearer-differenzengleichungen-ohne-anfangsbedingungen}

Zur Lösung der Differenzengleichung wird von der allgemeinen Form ausgegangen
\begin{align}
\sum_{n = 0}^{N}{c_{n}  y[ k - n]} &= \sum_{l = 0}^{L}{d_{l}  u[ k - l]} \label{eq:6.1}
\end{align}
Mit der Verschiebungsregel nach rechts ergibt sich die z-Transformierte der Gleichung zu
\begin{align}
\sum_{n = 0}^{N}{c_{n}  Y(z)  z^{- n}} &= \sum_{l = 0}^{L}{d_{l}  U(z)  z^{- l}} \label{eq:6.2}
\end{align}
Nach Ausklammern der beiden z-Transformierten \(Y(z)\) und \(U(z)\)
\begin{align}
Y(z)  \sum_{n = 0}^{N}{c_{n}  z^{- n}} &= U(z)  \sum_{l = 0}^{L}{d_{l}  z^{- l}} \label{eq:6.3}
\end{align}
kann die Gleichung nach der z-Transformierten des Ausgangssignals \(Y(z)\) aufgelöst werden.
\begin{align}
Y(z) &= \frac{\sum\limits_{l = 0}^{L}{d_{l}  z^{- l}}}{\sum_{n = 0}^{N}{c_{n}  z^{- n}}}  U(z) \label{eq:6.4}
\end{align}
Die Lösung im z-Bereich ist eine gebrochen rationale Funktion, die mit Hilfe der Partialbruchzerlegung und den bekannten Korrespondenzen in den Zeitbereich zurück transformiert werden kann. Dieses Vorgehen wird an der Differenzengleichung eines rekursiven Tiefpasses verdeutlicht.

\begin{example}{Rekursiver Tiefpass ohne Anfangsbedingungen}
\label{beispiel-rekursiver-tiefpass-ohne-anfangsbedingungen}

Es soll die Reaktion eines rekursiven Tiefpasses auf einen Sprung der Eingangsgröße \(u[k]\) von der Höhe \(U_{0}\) berechnet werden. Das Eingangssignal besitzt die z-Transformierte
\begin{align}
U(z) &= \frac{z}{z - 1}  U_{0} \label{eq:6.5}
\end{align}
Das Übertragungsverhalten des rekursiven Tiefpasses wird mit der Differenzengleichung
\begin{align}
y[ k] - GF  y[ k - 1] &= (1 - GF)  u[ k] \label{eq:6.6}
\end{align}
beschrieben. Die Transformation der Gleichung in den z-Bereich ergibt
\begin{align}
Y(z)  \left( 1 - GF  z^{- 1} \right) &= (1 - GF)  \frac{z}{z - 1}  U_{0} \label{eq:6.7}
\end{align}
Auflösen nach der z-Transformierten des Ausgangssignals \(Y(z)\) führt zu
\begin{align}
Y(z) &= \frac{z}{z - GF}  (1 - GF)  \frac{z}{z - 1}  U_{0} \label{eq:6.8}
\end{align}
Mit Hilfe der Partialbruchzerlegung kann \(Y(z)\) dargestellt werden als
\begin{align}
Y(z) &= \frac{z}{z - GF}  (1 - GF)  \frac{z}{z - 1}  U_{0} = \left( ( - GF)  \frac{z}{z - GF} + \frac{z}{z - 1} \right)  U_{0} \label{eq:6.9}
\end{align}
Die Darstellung mit Partialbrüchen ermöglicht die Rücktransformation in den Zeitbereich zu
\begin{align}
y[ k] &= \left( ( - GF)  GF^{k}  \sigma[ k] + \sigma[ k] \right)  U_{0} = \left( 1 - GF^{k + 1} \right)  U_{0}  \sigma[ k] \label{eq:6.10}
\end{align}
Ein Vergleich mit der über die Faltungssumme berechneten Systemantwort in Bild \ref{fig:6.2} zeigt identische Ausgangssignale.

\manualFigure
  [alt={Ausgangssignal eines rekursiven Tiefpasses mit GF = 0.8 bei einem Eingangssprung (\(U_{0} = 1\))}]
  {./media/image114.png}
  {Ausgangssignal eines rekursiven Tiefpasses mit GF = 0.8 bei einem Eingangssprung (\(U_{0} = 1\))}
  {\label{fig:6.2}}


\end{example}

\subsection{Lösung linearer Differenzengleichungen mit Anfangsbedingungen}\label{sec:luxf6sung-linearer-differenzengleichungen-mit-anfangsbedingungen}

Zur Lösung linearer Differenzengleichungen mit Anfangsbedingungen wird wieder von der allgemeinen Form der Differenzengleichung ausgegangen
\begin{align}
\sum_{n = 0}^{N}{c_{n}  y[ k - n]} &= \sum_{l = 0}^{L}{d_{l}  u[ k - l]} \label{eq:6.11}
\end{align}
Aufgrund der Zeitinvarianz kann das System um N nach links verschoben werden. Es ergibt sich die Gleichung
\begin{align}
\sum_{n = 0}^{N}{c_{n}  y[ k - n + N]} &= \sum_{l = 0}^{L}{d_{l}  u[ k - l + N]} \label{eq:6.12}
\end{align}
Da das Argument der Folge \(y[k -- n + N]\) damit immer größer gleich k ist, muss bei der Transformation der Gleichung in den z-Bereich die Verschiebungsregel nach links angewendet werden. Im z-Bereich ergibt sich die Gleichung
\begin{align}
\sum_{n = 0}^{N}{c_{n}  \left( Y(z)  z^{N - n} - z^{N - n}  \sum_{k = 0}^{N - n - 1}{y[ k]  z^{- k}} \right)} &= \sum_{l = 0}^{L}{d_{l}  \left( U(z)  z^{N - l} - z^{N - l}  \sum_{k = 0}^{N - l - 1}{u[ k]  z^{- k}} \right)} \label{eq:6.13}
\end{align}
Dieses Vorgehen ähnelt sehr der Lösung von Differentialgleichungen mit Anfangsbedingungen bei der Laplace-Transformation. In der allgemeinen Form erscheint es aufwendig. Bei praktischer Anwendung bestehen die Summen mit den Anfangswerten aber oft nur aus einzelnen Werten, sodass sich die Gleichung stark vereinfacht.

\begin{example}{RC-Tiefpass mit Anfangsbedingungen}
\label{beispiel-rc-tiefpass-mit-anfangsbedingungen}

Das Vorgehen wird wieder an der Differenzengleichung eines rekursiven Tiefpasses verdeutlicht.
\begin{align}
y[ k] - GF  y[ k - 1] &= (1 - GF)  u[ k] \label{eq:6.14}
\end{align}
Es soll die Reaktion des Tiefpasses auf einen Sprung der Eingangsgröße \(u[k]\) von der Höhe \(U_{0}\) und für die Anfangsbedingungen \(y[0]\) = 2 berechnet werden. Zur Berücksichtigung der Anfangsbedingungen muss bei der Transformation der Gleichung in den z-Bereich die Verschiebungsregel nach links angewendet werden. Eine Index-Transformation führt zu
\begin{align}
y[ k + 1] - GF  y[ k] &= (1 - GF)  u[ k + 1] \label{eq:6.15}
\end{align}
Die Transformation in den z-Bereich ergibt
\begin{align}
z  Y(z) - z  y[ 0] - GF  Y(z) &= (1 - GF)  \left( z  U(z) - z  u[ 0] \right) \label{eq:6.16}
\end{align}
Durch Einsetzen von \(U(z)\) und \(u[0]\) = 0~sowie Umformung
\begin{align}
(z - GF)  Y(z) &= (1 - GF)  \left( \frac{z^{2}}{z - 1}  U_{0} - z  U_{0} \right) + z  y[ 0] = (1 - GF)  U_{0}  \left( \frac{z^{2}}{z - 1} \right) + z  y[ 0] \label{eq:6.17}
\end{align}
ergibt sich die Lösung im z-Bereich
\begin{align}
Y(z) &= (1 - GF)  U_{0}  \left( \frac{z^{2}}{(z - 1)  (z - GF)} \right) + \frac{z}{z - GF}  y[ 0] \label{eq:6.18}
\end{align}
Der erste Term der z-Transformierten ist identisch zu dem Fall ohne Anfangsbedingungen in Gleichung (\ref{eq:6.8}). Der zweite Teil beschreibt das Systemverhalten aufgrund der Anfangsbedingung und kann mit der Korrespondenztafel in den Zeitbereich zurücktransformiert werden. Es ergibt sich das Ausgangssignal
\begin{align}
y[ k] &= \left( 1 - GF^{k + 1} \right)  U_{0}  \sigma[ k] + y[ 0]  GF^{k}  \sigma[ k] \label{eq:6.19}
\end{align}
Die Anfangsbedingung klingt ähnlich wie im zeitkontinuierlichen Bereich exponentiell ab. Bild \ref{fig:6.3} stellt die beiden Signalanteile dar.

\manualFigure
  [alt={Ausgangssignal eines rekursiven Tiefpasses mit GF = 0.8\\
bei einem Eingangssprung und Anfangsbedingung \(y[0]\) = 2}]
  {./media/image115.png}
  {Ausgangssignal eines rekursiven Tiefpasses mit GF = 0.8\\
bei einem Eingangssprung und Anfangsbedingung \(y[0]\) = 2}
  {\label{fig:6.3}}


\end{example}

\subsection{\texorpdfstring{Vorgehen zur Lösung linearer Differenzengleichungen }{Vorgehen zur Lösung linearer Differenzengleichungen }}\label{vorgehen-zur-luxf6sung-linearer-differenzengleichungen}

Die beiden Abschnitte zeigen das Vorgehen zur Lösung linearer Differenzengleichungen ohne und mit Anfangsbedingungen, das in Tabelle \ref{tab:6.1} zusammengefasst ist.

\begin{table}
  \centering
  \caption{Vorgehen bei der Berechnung der Systemantwort mit der z-Transformation}
  \label{tab:6.1}
  \begin{tabular}{lll}
    \toprule
    Schritt & \multicolumn{2}{l}{%
Lösung einer Differenzengleichung} \\
    \midrule
    & ohne Anfangsbedingungen & mit Anfangsbedingungen \\
    1 & Differenzengleichung & Verschobene Differenzengleichung \\
    & \(\sum_{n = 0}^{N}{c_{n}  y\lbrack k - n\rbrack} = \sum_{l = 0}^{L}{d_{l}  u\lbrack k - l\rbrack}\) & \(\sum_{n = 0}^{N}{c_{n}  y\lbrack k - n + N\rbrack} = \sum_{l = 0}^{L}{d_{l}  u\lbrack k - l + N\rbrack}\) \\
    2 & \multicolumn{2}{l}{%
Transformation der Differenzengleichung in den z-Bereich unter Verwendung der entsprechenden Verschiebungsregel} \\
    3 & \multicolumn{2}{l}{%
Berechnung der z-Transformierten \(Y(z)\)} \\
    4 & \multicolumn{2}{l}{%
Rücktransformation mit Hilfe der Partialbruchzerlegung und den bekannten Korrespondenzen} \\
    \bottomrule
  \end{tabular}
\end{table}

\section{Übertragungsfunktion zeitdiskreter Systeme}\label{sec:6.2}

Im zeitkontinuierlichen Bereich werden Systeme durch die Übertragungsfunktion \(G(s)\) charakterisiert. Sie ergibt sich aus der Laplace-Transformierten der Differentialgleichung bei verschwindenden Anfangsbedingungen. Dasselbe Verfahren wird auch bei zeitdiskreten Systemen angewendet. Zur Berechnung der z-Transformierten \(Y(z)\) kann die Verschiebungsregel verwendet werden. Im Fall verschwindender Anfangsbedingungen ergibt sich aus der Differenzengleichung
\begin{align}
\sum_{n = 0}^{N}{c_{n}  y[ k - n]} &= \sum_{l = 0}^{L}{d_{l}  u[ k - l]} \label{eq:6.20}
\end{align}
die Darstellung im z-Bereich zu
\begin{align}
Y(z)  \sum_{n = 0}^{N}{c_{n}  z^{- n}} &= U(z)  \sum_{l = 0}^{L}{d_{l}  z^{- l}} \label{eq:6.21}
\end{align}
Alternativ kann die Berechnung über die Faltungsregel durchgeführt werden. Aus der Faltung im Zeitbereich
\begin{align}
y[ k] &= \sum_{\kappa = 0}^{k}{g[\kappa]  u[ k - \kappa]} = \sum_{\kappa = 0}^{k}{g[ k - \kappa]  u[\kappa]} \label{eq:6.22}
\end{align}
folgt im z-Bereich
\begin{align}
Y(z) &= G(z)  U(z) \label{eq:6.23}
\end{align}
Durch einen Vergleich der beiden Darstellungen ergibt sich die Übertragungsfunktion für zeitdiskrete Systeme aus der Gleichung
\begin{align}
Y(z) &= G(z)  U(z) = \frac{\sum_{l = 0}^{L}{d_{l}  z^{- l}}}{\sum_{n = 0}^{N}{c_{n}  z^{- n}}}  U(z) \label{eq:6.24}
\end{align}
Die Übertragungsfunktion beschreibt den Zusammenhang zwischen Ein- und Ausgangssignal im z-Bereich.

\subsection{Bestimmung der Differenzengleichung aus der Übertragungsfunktion}\label{sec:bestimmung-der-differenzengleichung-aus-der-uxfcbertragungsfunktion}

Wie später gezeigt wird, ist das Ergebnis eines Filterentwurfs eine Übertragungsfunktion \(G(z)\), die die definierten Filtereigenschaften besitzt. Um den Filter zum Beispiel in Form eines Mikro-Controller-Programms realisieren zu können, muss die Übertragungsfunktion
\begin{align}
G(z) &= \frac{Y(z)}{U(z)} = \frac{\sum_{l = 0}^{L}{d_{l}  z^{- l}}}{\sum_{n = 0}^{N}{c_{n}  z^{- n}}} \label{eq:6.25}
\end{align}
in eine Differenzengleichung umgeformt werden. Durch Ausmultiplizieren der beiden Brüche ergibt sich die Gleichung
\begin{align}
Y(z)  \sum_{n = 0}^{N}{c_{n}  z^{- n}} &= U(z)  \sum_{l = 0}^{L}{d_{l}  z^{- l}} \label{eq:6.26}
\end{align}
Sie kann mit der Verschiebungsregel in den Zeitbereich zurücktransformiert werden.
\begin{align}
\sum_{n = 0}^{N}{c_{n}  y[ k - n]} &= \sum_{l = 0}^{L}{d_{l}  u[ k - l]} \label{eq:6.27}
\end{align}
Es wird deutlich, dass die Koeffizienten von Differenzengleichung und Übertragungsfunktion identisch sind. Zur Realisierung eines Filters wird der aktuelle Ausgangswert aus den vergangenen Ein- und Ausgangswerten berechnet. Dabei kann ohne Einschränkung der Allgemeinheit \(c_{0} = 1\) gesetzt werden.
\begin{align}
y[ k] &= \sum_{l = 0}^{L}{d_{l}  u[ k - l]} - \sum_{n = 1}^{N}{c_{n}  y[ k - n]} \label{eq:6.28}
\end{align}
Die Gleichung ist Grundlage für die technische Realisierung zeitdiskreter Systeme.

\begin{example}{Systemrealisierung}
\label{beispiel-systemrealisierung}

Gegeben ist ein System mit der Übertragungsfunktion
\begin{align}
G(z) &= \frac{Y(z)}{U(z)} = \frac{0.5  z}{z - 0.5} = \frac{0.5}{1 - 0.5  z^{- 1}} \label{eq:6.29}
\end{align}
Ausmultiplizieren
\begin{align}
Y(z)  \left( 1 - 0.5  z^{- 1} \right) &= 0.5  U(z) \label{eq:6.30}
\end{align}
und Rücktransformation ergibt mit der Verschiebungsregel
\begin{align}
y[ k] - 0.5  y[ k - 1] &= 0.5  u[ k] \label{eq:6.31}
\end{align}
Zur Realisierung des Filters als Algorithmus muss die Gleichung nach \(y[k]\) aufgelöst werden. Das aktuelle Ausgangssignal des Filters errechnet sich damit zu
\begin{align}
y[ k] &= \frac{1}{2}  \left( u[ k] + y[ k - 1] \right) \label{eq:6.32}
\end{align}
Diese Gleichung kann zum Beispiel als Software in einem Controller oder als zeitdiskrete Schaltung realisiert werden.

\end{example}

\subsection{Impuls- und Sprungantwort}\label{sec:impuls--und-sprungantwort}

Die Übertragungsfunktion \(G(z)\) ist die z-Transformierte der Impulsantwort. Das ergibt sich unmittelbar aus der Definitionsgleichung, wenn für \(U(z)\) die z-Transformierte des Impulses
\begin{align}
U(z) &= 1 \label{eq:6.33}
\end{align}
eingesetzt wird:
\begin{align}
Y(z) &= G(z)  U(z) = G(z)  1 = G(z) \label{eq:6.34}
\end{align}
Bei Anregung des Systems mit einer Sprungfolge mit der z-Transformierten
\begin{align}
U(z) &= \frac{z}{z - 1} \label{eq:6.35}
\end{align}
ergibt sich die z-Transformierte des Ausgangssignals zu
\begin{align}
Y(z) &= G(z)  U(z) = G(z)  \frac{z}{z - 1} = H(z) \label{eq:6.36}
\end{align}
Diese Gleichung kann mit der Summationsregel rücktransformiert werden, und es ergibt sich
\begin{align}
h[ k] &= \sum_{\kappa = 0}^{k}{g[\kappa]} \label{eq:6.37}
\end{align}
Die Sprungantwort ergibt sich aus der Summe der einzelnen Folgenwerte der Impulsantwort. Dieses Ergebnis entspricht sinngemäß den Eigenschaften zeitkontinuierlicher Systeme. Dort ist die Sprungantwort \(h(t)\) das Integral der Impulsantwort \(g(t)\). Umgekehrt kann die Impulsantwort aus der Sprungantwort berechnet werden. Im z-Bereich gilt:
\begin{align}
H(z) &= \frac{z}{z - 1}  G(z) = \frac{1}{1 - z^{- 1}}  G(z) \label{eq:6.38}
\end{align}
Diese Gleichung kann in den Zeitbereich zurücktransformiert werden.
\begin{align}
g[ k] &= h[ k] - h[ k - 1] \label{eq:6.39}
\end{align}
Auch dieses Ergebnis entspricht sinngemäß den Eigenschaften zeitkontinuierlicher Systeme. Dort ist die Impulsantwort \(g(t)\) die erweiterte Ableitung der Sprungantwort \(h(t)\). Für zeitkontinuierliche Systeme wird der Endwert der Sprungantwort als Verstärkung definiert. Um den Begriff der Verstärkung für zeitdiskrete Systeme zu definieren, wird der Grenzwert untersucht.
\begin{align}
\lim_{k \rightarrow \infty}h[ k] &= \lim_{z \rightarrow 1}(z - 1)  H(z) = G(1) \label{eq:6.40}
\end{align}
Die Verstärkung eines zeitdiskreten Systems ergibt sich aus dem Wert der Übertragungsfunktion an der Stelle \(z = 1\). Tabelle \ref{tab:6.2} fasst die Zusammenhänge von Impuls- und Sprungantwort zusammen.

\begin{table}
  \centering
  \caption{Übersicht zum Zusammenhang von Impuls- und Sprungantwort}
  \label{tab:6.2}
  \begin{tabular}{lll}
    \toprule
    Eigenschaft & Zeitbereich & z-Bereich \\
    \midrule
    Impulsantwort & \(g\lbrack k\rbrack = h\lbrack k\rbrack - h\lbrack k - 1\rbrack\) & \(G(z) = H(z)  \left( 1 - z^{- 1} \right)\) \\
    Sprungantwort & \(h\lbrack k\rbrack = \sum_{\kappa = 0}^{k}{g\lbrack\kappa\rbrack}\) & \(H(z) = G(z)  \frac{z}{z - 1}\) \\
    Verstärkung des Systems & \(\lim_{k \rightarrow \infty}h\lbrack k\rbrack\) & \(\lim_{z \rightarrow 1}(z - 1)  H(z) = G(1)\) \\
    \bottomrule
  \end{tabular}
\end{table}

Zur Vertiefung des Begriffes der Sprungantwort werden ein zeitkontinuierliches und ein zeitdiskretes System verglichen, die an den Stellen t = kT\textsubscript{A} dieselben Werte aufweisen.
\begin{align}
g[ k] &= g\left( k  T_{A} \right) \label{eq:6.41}
\end{align}
Die Sprungantwort des zeitkontinuierlichen Systems ergibt sich aus dem Integral
\begin{align}
h(t) &= \int_{0}^{t}{g(\tau)}d\tau \label{eq:6.42}
\end{align}
Die Sprungantwort des zeitdiskreten Systems ergibt sich aus der Summe
\begin{align}
h[ k] &= \sum_{\kappa = 0}^{k}{g[\kappa]} \label{eq:6.43}
\end{align}
Während das Integral die Fläche unter der Impulsantwort bis zum Punkt t repräsentiert, nähert die Summe die Fläche über Rechtecke. Die Breite der Rechtecke ist bei zeitdiskreten Systemen auf den Wert eins normiert. Aus diesem Grund gilt:
\begin{align}
h[ k] \neq h\left( k  T_{A} \right) \label{eq:6.44}
\end{align}
Die Abtastzeit kann dadurch berücksichtigt werden, dass die Sprungantwort \(h[k]\) mit der Abtastzeit \(T_{A}\) multipliziert wird. Damit wird die Breite der Rechtecke von 1 auf \(T_{A}\) geändert.
\begin{align}
h[ k]  T_{A} \approx h\left( k  T_{A} \right) \label{eq:6.45}
\end{align}
Auch nach dieser Korrektur sind die Sprungantworten nicht identisch, da die Fläche unter der zeitkontinuierlichen Impulsantwort im zeitdiskreten Fall nur über Rechtecke approximiert wird.

\begin{example}{Vergleich von Sprungantworten}
\label{beispiel-vergleich-von-sprungantworten}

Gegeben ist ein zeitkontinuierliches System mit der Impulsantwort
\begin{align}
g(t) &= 2  e^{- \frac{}{}\frac{t}{10\,\text{ms}}}  \sigma(t) \label{eq:6.46}
\end{align}
Das System wird mit einer Abtastzeit \(T_{A} = 1\,\text{ms}\) abgetastet. Es ergeben sich die Folgenwerte
\begin{align}
g[ k] &= 2  e^{- \frac{k}{10}}  \sigma[ k] \label{eq:6.47}
\end{align}
Die Sprungantwort des zeitkontinuierlichen Systems lautet
\begin{align}
h(t) &= \int_{0}^{t}{2  e^{- \frac{\tau}{10\,\text{ms}}}d\tau} = - 20\,\text{ms}  \left. \ e^{- \frac{\tau}{10\,\text{ms}}} \right|_{0}^{t} = 20\,\text{ms}  \left( 1 - e^{- \frac{t}{10\,\text{ms}}} \right)  \sigma(t) \label{eq:6.48}
\end{align}
Die Sprungantwort des zeitdiskreten Systems errechnet sich zu
\begin{align}
h[ k] &= \sum_{\kappa = 0}^{k}{2  e^{- \frac{\kappa}{10}}} = 2  {\sum_{\kappa = 0}^{k}\left( e^{- \frac{1}{10}} \right)}^{\kappa} = 2  \frac{1 - e^{- \frac{k + 1}{10}}}{1 - e^{- \frac{1}{10}}} \label{eq:6.49}
\end{align}
Durch die Multiplikation der Sprungantwort mit der Abtastzeit \(T_{A}\) ergibt sich die Näherung
\begin{align}
h\left( k  T_{A} \right) \approx h[ k]  T_{A} &= 2\,\text{ms}  \frac{1 - e^{- \frac{k + 1}{10}}}{1 - e^{- \frac{1}{10}}} \label{eq:6.50}
\end{align}
Die Impuls- und Sprungantworten sind in Bild \ref{fig:6.4} dargestellt.

\manualFigure
  [alt={Vergleich der Sprungantworten eines zeitdiskreten und eines zeitkontinuierlichen Systems\\
mit identischen Impulsantworten \(g[k]\) = g(kT\textsubscript{A})}]
  {./media/image116.png}
  {Vergleich der Sprungantworten eines zeitdiskreten und eines zeitkontinuierlichen Systems\\
mit identischen Impulsantworten \(g[k]\) = g(kT\textsubscript{A})}
  {\label{fig:6.4}}


Es wird deutlich, dass die Sprungantworten nur näherungsweise gleich sind. Grund dafür ist die Approximation der Fläche unter der Impulsantwort mit den in Bild \ref{fig:6.4} eingezeichneten Rechtecken. Ein ausführlicher Vergleich zwischen zeitkontinuierlichen und zeitdiskreten Systemen wird in Kapitel \hyperref[_Ref411155749]{0} vorgenommen.

\end{example}

\section{Interpretation der Übertragungsfunktion}\label{sec:6.3}

Gebrochen rationale Übertragungsfunktionen charakterisieren lineare, zeitinvariante Systeme im z-Bereich. An der Übertragungsfunktion lassen sich auch ohne Rücktransformation in den Zeitbereich wichtige Systemeigenschaften ablesen. In diesem Abschnitt werden wichtige Zusammenhänge zwischen der Übertragungsfunktion im z-Bereich und dem Systemverhalten im Zeitbereich hergeleitet

\subsection{Exkurs in die Darstellungsformen von Übertragungsfunktionen}\label{sec:exkurs-in-die-darstellungsformen-von-uxfcbertragungsfunktionen}

Bei der Diskussion der Übertragungsfunktion werden unterschiedliche Darstellungsformen genutzt. Sie werden an dieser Stelle miteinander verglichen und ihre Bedeutung für die Systemtheorie herausgearbeitet.

\subsubsection{Darstellungsformen zeitdiskreter LTI-Systeme}\label{sec:darstellungsformen-zeitdiskreter-lti-systeme}

Die Übertragungsfunktion von zeitdiskreten LTI-Systemen kann auf unterschiedliche Art dargestellt werden. Die Beschreibung eines Systems mit der Differenzengleichung
\begin{align}
\sum_{n = 0}^{N}{c_{n}  y[ k - n]} &= \sum_{l = 0}^{L}{d_{l}  u[ k - l]} \label{eq:6.51}
\end{align}
ist direkt mit einer Übertragungsfunktion der Form
\begin{align}
G(z) &= \frac{Y(z)}{U(z)} = \frac{\sum_{l = 0}^{L}{d_{l}  z^{- l}}}{\sum_{n = 0}^{N}{c_{n}  z^{- n}}} \label{eq:6.52}
\end{align}
verknüpft. Um ein System mit Hilfe von Speichergliedern zu realisieren, wird die Gleichung nach \(Y(z)\) aufgelöst.
\begin{align}
Y(z) &= \sum_{l = 0}^{L}{d_{l}  U(z)  z^{- l}} - \sum_{n = 1}^{N}{c_{n}  Y(z)  z^{- n}} = d_{0}  U(z) + \sum_{n = 1}^{N}{\left( d_{n}  U(z) - c_{n}  Y(z) \right)  z^{- n}} \label{eq:6.53}
\end{align}
Es ergibt sich die in Bild \ref{fig:6.5} gezeigte kanonische Darstellungsform von Systemen mit Speichergliedern.

\begin{figure}
  \centering
  \includegraphics[alt={Kanonisches Blockschaltbild eines linearen, zeitinvarianten Systems im zeitkontinuierlichen Bereich}]{./media/6_2_kanonisches_blockschaltbild_eine.pdf}
  \caption{Kanonisches Blockschaltbild eines linearen, zeitinvarianten Systems im zeitkontinuierlichen Bereich}
  \label{fig:6.5}
\end{figure}

Durch Erweiterung der Übertragungsfunktion (\ref{eq:6.52}) mit \(z^{N}\) ergibt sich eine Übertragungsfunktion, bei der für \(d_{0}\) ≠ 0 Zähler- und Nennergrad identisch sind.
\begin{align}
G(z) &= \frac{Y(z)}{U(z)} = \frac{z^{N}  \sum_{l = 0}^{L}{d_{l}  z^{- l}}}{z^{N}  \sum_{n = 0}^{N}{c_{n}  z^{- n}}} = \frac{\sum_{m = 0}^{M}{b_{m}  z^{m}}}{\sum_{n = 0}^{N}{a_{n}  z^{n}}} \label{eq:6.54}
\end{align}
Diese Darstellungsform wird bei der Rücktransformation einer z-Transformierten \(G(z)\) durch Partialbruchzerlegung eingesetzt.

\begin{example}{Darstellungsformen von Übertragungsfunktionen}
\label{beispiel-darstellungsformen-von-uxfcbertragungsfunktionen}

Gegeben ist ein System mit der Differenzengleichung
\begin{align}
y[ k] + 3  y[ k - 1] + 2  y[ k - 2] &= u[ k] + 4  u[ k - 1] \label{eq:6.55}
\end{align}
und der daraus resultierenden Übertragungsfunktion
\begin{align}
G(z) &= \frac{Y(z)}{U(z)} = \frac{1 + 4  z^{- 1}}{1 + 3  z^{- 1} + 2  z^{- 2}} \label{eq:6.56}
\end{align}
Auflösen nach \(Y(z)\) ergibt
\begin{align}
Y(z) &= U(z) + \left( 4  U(z) - 3  Y(z) - 2  Y(z)  z^{- 1} \right)  z^{- 1} \label{eq:6.57}
\end{align}
Bild \ref{fig:6.6} zeigt die kanonische Darstellungsform des Systems mit Speichergliedern.

\manualFigure
  [alt={Kanonisches Blockschaltbild Systems mit der Differenzengleichung (\ref{eq:6.55})}]
  {./media/image118.pdf}
  {Kanonisches Blockschaltbild Systems mit der Differenzengleichung (\ref{eq:6.55})}
  {\label{fig:6.6}}


Die Übertragungsfunktion hat im Zähler eine Ordnung von \(M = 1\) und im Nenner eine Ordnung\\
\(N = 2\). Erweitern der Übertragungsfunktion mit z² führt zu der Darstellung
\begin{align}
G(z) &= \frac{z^{2} + 4  z}{z^{2} + 3  z + 2} \label{eq:6.58}
\end{align}
Zähler- und Nennerpolynom haben denselben Grad \(N = 2\). Nach Polynomdivision ergibt sich eine echt gebrochen rationale Funktion
\begin{align}
G(z) &= \frac{z^{2} + 4  z}{z^{2} + 3  z + 2} = 1 + \frac{z - 2}{z^{2} + 3  z + 2} \label{eq:6.59}
\end{align}
Sie kann als Summe von Partialbrüchen dargestellt
\begin{align}
G(z) &= 1 + \frac{z - 2}{z^{2} + 3  z + 2} = 1 - \frac{3}{z + 1} + \frac{4}{z + 2} = 1 - 3  \frac{z}{z + 1}  z^{- 1} + 4  \frac{z}{z + 2}  z^{- 1} \label{eq:6.60}
\end{align}
und in den Zeitbereich zurücktransformiert werden.
\begin{align}
g[ k] &= \delta[ k] - 3  ( - 1)^{k - 1}  \sigma[ k - 1] + 4  ( - 2)^{k - 1}  \sigma[ k - 1] \label{eq:6.61}
\end{align}
\end{example}

\subsubsection{Darstellungsformen zeitkontinuierlicher LTI-Systeme}\label{sec:darstellungsformen-zeitkontinuierlicher-lti-systeme}

Auch bei zeitkontinuierlichen Systemen kann die Übertragungsfunktion auf unterschiedliche Arten dargestellt werden. In Teil A dieser Buchreihe wird in Kapitel 5 ausgehend von einer Differentialgleichung N-ter Ordnung
\begin{align}
\sum_{n = 0}^{N}{a_{n}  \frac{d^{n}y}{dt^{n}}} &= \sum_{m = 0}^{M}{b_{m}  \frac{d^{m}u}{dt^{m}}} \label{eq:6.62}
\end{align}
die Übertragungsfunktion
\begin{align}
G(s) &= \frac{Y(s)}{U(s)} = \frac{\sum_{m = 0}^{M}{b_{m}  s^{m}}}{\sum_{n = 0}^{N}{a_{n}  s^{n}}} \label{eq:6.63}
\end{align}
bestimmt. Die Interpretation des Systems erfolgt anhand dieser Darstellungsform mit Zähler- und Nennerpolynom. Da die Differentiation im Zeitbereich nicht kausal ist, kann diese Darstellungsform jedoch nicht zur Simulation zeitkontinuierlicher LTI-Systeme verwendet werden. Deshalb wird im Abschnitt 3.5 von Teil A dieser Buchreihe die Differentialgleichung N-mal integriert. Analog zu Gleichung (\ref{eq:6.52}) ergibt sich der Ausdruck
\begin{align}
G(s) &= \frac{Y(s)}{U(s)} = \frac{\sum_{m = 0}^{M}{b_{m}  s^{m}}}{\sum_{n = 0}^{N}{a_{n}  s^{n}}} = \frac{\frac{1}{s^{N}}  \sum_{m = 0}^{M}{b_{m}  s^{m}}}{\frac{1}{s^{N}}  \sum_{n = 0}^{N}{a_{n}  s^{n}}} = \frac{\sum_{m = 0}^{M}{b_{m}  s^{m - N}}}{\sum_{n = 0}^{N}{a_{n}  s^{n - N}}} = \frac{\sum_{l = 0}^{N}{d_{l}  s^{- l}}}{\sum_{n = 0}^{N}{c_{n}  s^{- n}}} \label{eq:6.64}
\end{align}
Ausmultiplizieren der Gleichung
\begin{align}
Y(s)  \sum_{n = 0}^{N}{c_{n}  s^{- n}} &= U(s)  \sum_{l = 0}^{N}{d_{l}  s^{- l}} \label{eq:6.65}
\end{align}
und Auflösen nach \(Y(s)\) führt mit \(c_{0} = 1\) zu dem Ausdruck
\begin{align}
{Y\left(s\right)} &=  {U\left(s\right) \sum _{l=0}^{N}d_{l}  s^{-l}  -Y\left(s\right) \sum _{n=1}^{N}c_{n}  s^{-n}  } \\
{} &=  {\sum _{l=0}^{N}d_{l}  U\left(s\right) s^{-l}  -\sum _{n=1}^{N}c_{n}  Y\left(s\right) s^{-n}  }
 \label{eq:6.66}
\end{align}
Um die Integrierer zusammenfassen zu können, wird die erste Summe zerlegt.
\begin{align}
Y(s) &= \sum_{l = 0}^{N}{d_{l}  U(s)  s^{- l}} - \sum_{n = 1}^{N}{c_{n}  Y(s)  s^{- n}} = d_{0}  U(s) + \sum_{l = 1}^{N}{d_{l}  U(s)  s^{- l}} - \sum_{n = 1}^{N}{c_{n}  Y(s)  s^{- n}} \label{eq:6.67}
\end{align}
Damit können die beiden übrigen Summen zusammengefasst werden.
\begin{align}
Y(s) &= d_{0}  U(s) + \sum_{n = 1}^{N}{\left( d_{l}  U(s) - c_{n}  Y(s) \right)  s^{- n}} \label{eq:6.68}
\end{align}
Es ergibt sich die in Bild \ref{fig:6.7} gezeigte kanonische Darstellungsform von Systemen mit Integrierern.

\begin{figure}
  \centering
  \includegraphics[alt={Kanonisches Blockschaltbild eines linearen, zeitinvarianten Systems im zeitkontinuierlichen Bereich}]{./media/6_3_kanonisches_blockschaltbild_eine.pdf}
  \caption{Kanonisches Blockschaltbild eines linearen, zeitinvarianten Systems im zeitkontinuierlichen Bereich}
  \label{fig:6.7}
\end{figure}

\begin{example}{Darstellungsformen der Übertragungsfunktion für ein Feder-Masse-Dämpfer-Systems}
\label{beispiel-darstellungsformen-der-uxfcbertragungsfunktion-fuxfcr-ein-feder-masse-duxe4mpfer-systems}

Die unterschiedlichen Darstellungsformen für zeitkontinuierliche Systeme werden anhand eines Feder-Masse-Dämpfer-Systems mit der Differentialgleichung
\begin{align}
F_{E}(t) &= m  \frac{d^{2}x}{dt^{2}} + D  \frac{dx}{dt} + c  x(t) \label{eq:6.69}
\end{align}
verdeutlicht. Eingangssignal ist der Kraftverlauf \(F_{E}(t)\), Ausgangssignal ist die Auslenkung \(x(t)\). Transformation in den Laplace-Bereich ergibt bei verschwindenden Anfangsbedingungen die Gleichung
\begin{align}
F_{E}(s) &= m  s^{2}  X(s) + D  s  X(s) + c  X(s) \label{eq:6.70}
\end{align}
Damit lautet die Übertragungsfunktion
\begin{align}
\frac{X(s)}{F_{E}(s)} &= \frac{1}{m  s^{2} + D  s + c} = \frac{1}{m}  \frac{1}{s^{2} + \frac{D}{m}  s + \frac{c}{m}} \label{eq:6.71}
\end{align}
Auf Basis dieser Darstellung und den Polen des Systems
\begin{align}
\alpha_{1,2} &= - \frac{D}{2\,\text{m}} \pm \sqrt{\frac{D^{2}}{4\,\text{m}^{2}} - \frac{c}{m}} \label{eq:6.72}
\end{align}
können die Systemeigenschaften diskutiert werden. Andererseits erfolgt die Systemsimulation über eine Darstellungsform
\begin{align}
\frac{X(s)}{F_{E}(s)} &= \frac{\frac{1}{s^{2}}}{m + \frac{D}{s} + \frac{c}{s^{2}}} \label{eq:6.73}
\end{align}
Sie kann nach \(X(s)\) aufgelöst werden
\begin{align}
X(s) &= \frac{1}{m}  \left( \frac{1}{s^{2}}  F_{E}(s) - \frac{D}{s}  X(s) - \frac{c}{s^{2}}  X(s) \right) \label{eq:6.74}
\end{align}
Es ergibt sich das in Bild \ref{fig:6.8} dargestellte kanonische Blockschaltbild.

\manualFigure
  [alt={Kanonisches Blockschaltbild eines Feder-Masse-Dämpfer-Systems}]
  {./media/image121.pdf}
  {Kanonisches Blockschaltbild eines Feder-Masse-Dämpfer-Systems}
  {\label{fig:6.8}}


\end{example}

Je nach Anwendungsfall wird die Darstellung mit Zähler- und Nennerpolynom oder die Darstellung über Integrierglieder mit der Laplace-Transformierten 1/s verwendet.

\subsubsection{Zusammenfassung der Darstellungsformen von Übertragungsfunktionen}\label{sec:zusammenfassung-der-darstellungsformen-von-uxfcbertragungsfunktionen}

Übertragungsfunktionen werden damit sowohl bei zeitkontinuierlichen als auch bei zeitdiskreten Systemen in zwei Darstellungsformen behandelt. Sie sind in Tabelle \ref{tab:6.3} zusammen mit ihrer Anwendung dargestellt.

\begin{table}
  \centering
  \caption{Übersicht zu Darstellungsformen von Übertragungsfunktionen und ihrer Anwendung}
  \label{tab:6.3}
  \begin{tabular}{lll}
    \toprule
    Anwendung & Zeitkontinuierliche Systeme & Zeitdiskrete Systeme \\
    \midrule
    Mathematische Diskussion der Übertragungsfunktion & \(G(s) = \frac{\sum_{m = 0}^{M}{b_{m}  s^{m}}}{\sum_{n = 0}^{N}{a_{n}  s^{n}}}\) & \(G(z) = \frac{\sum_{m = 0}^{M}{b_{m}  z^{m}}}{\sum_{n = 0}^{N}{a_{n}  z^{n}}}\) \\
    Darstellungsform zur 
Realisierung der Systeme
 & \(G(s) = \frac{\sum_{l = 0}^{L}{d_{l}  s^{- l}}}{\sum_{n = 0}^{N}{c_{n}  s^{- n}}}\) & \(G(z) = \frac{\sum_{l = 0}^{L}{d_{l}  z^{- l}}}{\sum_{n = 0}^{N}{c_{n}  z^{- n}}}\) \\
    \bottomrule
  \end{tabular}
\end{table}

Die folgende Diskussion der Übertragungsfunktion zeitdiskreter Systeme bezieht sich damit auf eine Übertragungsfunktion der Form
\begin{align}
G(z) &= \frac{\sum_{m = 0}^{M}{b_{m}  z^{m}}}{\sum_{n = 0}^{N}{a_{n}  z^{n}}} \label{eq:6.75}
\end{align}
\subsection{Pol-Nullstellen-Diagramme}\label{sec:pol-nullstellen-diagramme}

Die Interpretation der Übertragungsfunktion erlaubt Aussagen über Stabilität, Schwingungsneigung und Sprungfähigkeit von Systemen. Ein wesentlicher Schritt zur Interpretation des Systems ist die Analyse der unterschiedlichen Pole \(α_{n}\) und der unterschiedlichen Nullstellen \(β_{m}\). Dazu wird die Übertragungsfunktion durch Linearfaktoren dargestellt.
\begin{align}
G(z) &= \frac{B(z)}{A(z)} = \frac{\sum_{m = 0}^{M}{b_{m}  z^{m}}}{\sum_{n = 0}^{N}{a_{n}  z^{n}}} = \frac{b_{M}}{a_{N}}  \frac{\left( z - \beta_{1} \right)^{M_{1}}  \left( z - \beta_{2} \right)^{M_{2}}  ...}{\left( z - \alpha_{1} \right)^{N_{1}}  \left( z - \alpha_{2} \right)^{N_{2}}  ...} \label{eq:6.76}
\end{align}
Die Pole \(α_{n}\) und Nullstellen \(β_{m}\) können zur besseren Übersicht in der komplexen Ebene dargestellt werden. Dabei werden Nullstellen mit einem Kreis und Pole mit einem Kreuz dargestellt, ihre Vielfachheit \(N_{n}\) und \(M_{m}\) wird in Klammern angegeben. Die Diagramme werden als Pol-Nullstellen-Diagramme oder als Pole-Zero-Maps bezeichnet.

\begin{example}{Pol-Nullstellen-Diagramm}
\label{beispiel-pol-nullstellen-diagramm}

Ein System mit der Differenzengleichung
\begin{align}
y[ k] - 0.2  y[ k - 1] - 0.15  y[ k - 2] &= u[ k] + 1.5  u[ k - 1] \label{eq:6.77}
\end{align}
weist eine Übertragungsfunktion
\begin{align}
G(z) &= \frac{Y(z)}{U(z)} = \frac{1 + 1.5  z^{- 1}}{1 - 0.2  z^{- 1} - 0.15  z^{- 2}} = \frac{z^{2} + 1.5  z}{z^{2} - 0.2  z - 0.15} = \frac{z  (z + 1.5)}{(z - 0.5)  (z + 0.3)} \label{eq:6.78}
\end{align}
auf. Sie besitzt die Nullstellen \(β_{1}\) = - 1.5 und \(β_{2} = 0\), sowie die reellen Polstellen \(α_{1}\) = - 0.3 und\\
\(α_{2} = 0.5\). Es ergibt sich das in Bild \ref{fig:6.9} dargestellte Pol-Nullstellen-Diagramm.

\manualFigure
  [alt={Pol-Nullstellen-Diagramm des Systems aus Gleichung (\ref{eq:6.77})}]
  {./media/image122.png}
  {Pol-Nullstellen-Diagramm des Systems aus Gleichung (\ref{eq:6.77})}
  {\label{fig:6.9}}


Das Pol-Nullstellen-Diagramm bietet einen guten Überblick über die Lage der Pole und Nullstellen in der komplexen Ebene.

\end{example}

\subsection{Übertragungsfunktion und Kausalität von Systemen}\label{sec:uxfcbertragungsfunktion-und-kausalituxe4t-von-systemen}

Zur Überprüfung der Kausalität eines Systems mit der Übertragungsfunktion
\begin{align}
G(z) &= \frac{Y(z)}{U(z)} = \frac{\sum_{m = 0}^{M}{b_{m}  z^{m}}}{\sum_{n = 0}^{N}{a_{n}  z^{n}}} = \frac{\sum_{m = 0}^{M}{b_{m}  z^{m - N}}}{\sum_{n = 0}^{N}{a_{n}  z^{n - N}}} \label{eq:6.79}
\end{align}
werden Zähler- und Nennerpolynom durch \(z^{N}\) dividiert. Ausmultiplizieren und Auflösen der Summen führt zur Gleichung
\begin{align}
Y(z)  \left( a_{N} + ... + a_{0}  z^{- N} \right) &= U(z)  \left( b_{M}  z^{M - N} + ... + b_{0}  z^{- N} \right) \label{eq:6.80}
\end{align}
Rücktransformation ergibt für verschwindende Anfangsbedingungen
\begin{align}
a_{N}  y[ k] + ... + a_{0}  y[ k - N] &= b_{M}  u[ k + M - N] + ... + b_{0}  u[ k - N] \label{eq:6.81}
\end{align}
Das System ist kausal, wenn \(y[k]\) nur von aktuellen und vergangenen Eingangswerten abhängt. Das ist der Fall, wenn die Bedingung M ≤ N erfüllt ist. Der Zählergrad der Übertragungsfunktion darf maximal so groß sein wie der Nennergrad, damit das System kausal ist.

\begin{example}{Nicht kausales System}
\label{beispiel-nicht-kausales-system}

Gegeben ist ein System mit der Übertragungsfunktion
\begin{align}
G(z) &= \frac{Y(z)}{U(z)} = \frac{3  z^{2} + z + 2}{2  z + 1} \label{eq:6.82}
\end{align}
Der Zählergrad \(M = 2\) ist größer als der Nennergrad \(N = 1\). Das System ist demnach nicht kausal. Umformen der Übertragungsfunktion führt zu
\begin{align}
Y(z)  (2  z + 1) &= U(z)  \left( 3  z^{2} + z + 2 \right) \label{eq:6.83}
\end{align}
beziehungsweise
\begin{align}
Y(z)  \left( 2 + z^{- 1} \right) &= U(z)  \left( 3  z + 1 + 2  z^{- 1} \right) \label{eq:6.84}
\end{align}
Rücktransformation führt bei verschwindenden Anfangsbedingungen zu der Differenzengleichung
\begin{align}
2  y[ k] + y[ k - 1] &= 3  u[ k + 1] + u[ k] + 2  u[ k - 1] \label{eq:6.85}
\end{align}
Der aktuelle Ausgangswert
\begin{align}
y[ k] &= \frac{1}{2}  \left( 3  u[ k + 1] + u[ k] + 2  u[ k - 1] - y[ k - 1] \right) \label{eq:6.86}
\end{align}
ist von dem zukünftigen Wert \(u[k+1]\) abhängig, das System ist damit nicht kausal.

\end{example}

\subsection{Übertragungsfunktion mit Zählergrad M gleich Nennergrad N}\label{sec:uxfcbertragungsfunktion-mit-zuxe4hlergrad-m-gleich-nennergrad-n}

Für kausale Systeme ist der Zählergrad M der Übertragungsfunktion kleiner oder gleich dem Nennergrad N. Für den Fall \(M = N\) muss vor der Partialbruchzerlegung eine Polynomdivision durchgeführt werden. Dadurch entsteht im Ausdruck für die Übertragungsfunktion \(G(z)\) ein konstanter Summand. Da die inverse z-Transformierte zu einer Konstanten die Impulsfolge \(δ[k]\) ist, entspricht diesem Summand im Zeitbereich der Folgenwert für k = 0
\begin{align}
g[ k &= 0] = \frac{b_{M}}{a_{N}} \label{eq:6.87}
\end{align}
Alle weiteren Summanden der Partialbruchzerlegung werden im Zeitbereich mit Sprungfunktionen der Form \(σ[k - 1]\) multipliziert und sind erst für k \textgreater{} 0 von null verschieden. Damit ist ein System nur dann sprungfähig, wenn der Zählergrad M gleich dem Nennergrad N ist.

\begin{example}{Sprungfähiges System}
\label{beispiel-sprungfuxe4higes-system}

Die Übertragungsfunktion \(G(z)\) soll interpretiert werden. Da der Zählergrad genauso groß ist wie der Nennergrad, wird eine Polynomdivision durchgeführt. Um auf einen Ausdruck der Korrespondenztafel zu kommen, muss der Restbruch mit z erweitert werden.
\begin{align}
G(z) &= \frac{2 + z^{- 1}}{1 - 0.2  z^{- 1}} = \frac{2  z + 1}{z - 0.2} = 2 + \frac{1.4}{z - 0.2} = 2 + \frac{1.4  z}{z - 0.2}  z^{- 1} \label{eq:6.88}
\end{align}
Damit kann die Impulsantwort mit der Korrespondenztafel und der Verschiebungsregel bestimmt werden zu
\begin{align}
g[ k] &= 2  \delta[ k] + 1.4  0.2^{k - 1}  \sigma[ k - 1] \label{eq:6.89}
\end{align}
Sie springt an der Stelle \(k = 0\) auf den Wert 2.

\end{example}

\subsection{Partialbruchzerlegung mit einfachen reellen Polen}\label{sec:partialbruchzerlegung-mit-einfachen-reellen-polen}

Besitzt die Übertragungsfunktion einfache reelle Pole \(α_{n}\) ≠ 0, kann sie über die Partialbruchzerlegung dargestellt werden als
\begin{align}
G(z) &= \sum_{n = 1}^{N}\frac{A_{n}}{z - \alpha_{n}} \label{eq:6.90}
\end{align}
Jeder einzelne Partialbruch hat die Form
\begin{align}
G_{n}(z) &= \frac{A_{n}}{z - \alpha_{n}} = A_{n}  \frac{z}{z - \alpha_{n}}  z^{- 1} \label{eq:6.91}
\end{align}
Im Zeitbereich ergibt sich damit für jeden Partialbruch eine Exponentialfunktion, die um einen Folgenindex nach rechts verschoben ist:
\begin{align}
g_{n}[ k] &= A_{n}  \alpha_{n}^{k - 1}  \sigma[ k - 1] \label{eq:6.92}
\end{align}
Die Summe der einzelnen Partialbrüche entspricht deshalb im Zeitbereich der Folge
\begin{align}
g[ k] &= \sum_{n = 1}^{N}{A_{n}  \alpha_{n}^{k - 1}  \sigma[ k - 1]} \label{eq:6.93}
\end{align}
Der Betrag der Pole \(α_{n}\) entscheidet über die Frage, ob die Impulsantwort konvergent oder divergent ist und damit über die Stabilität des Systems.

Die Impulsantwort beginnt bei sprungfähigen Systemen bei \(k = 0\) und bei nicht sprungfähigen Systemen bei \(k = 1\), und sie ist unendlich lang. Deshalb wird die Impulsantwort als Infinite-Impulse-Response und die Systeme als Infinite-Impulse-Response-Systeme (IIR-Systeme) bezeichnet. IIR-Systeme sind daran zu erkennen, dass zumindest ein Pol \(α_{n}\) der Übertragungsfunktion nicht im Koordinatenursprung liegt.

\begin{example}{Übertragungsfunktion mit einfachen reellen Polen}
\label{beispiel-uxfcbertragungsfunktion-mit-einfachen-reellen-polen}

Die untenstehende Übertragungsfunktion \(G(z)\) soll interpretiert werden. Ihr Zählergrad ist kleiner als der Nennergrad, und sie hat zwei einfache reelle Pole. Die Übertragungsfunktion kann in Partialbrüche zerlegt werden.
\begin{align}
{G(z) = \frac{3  z^{- 1} - 3.5  z^{- 2}}{1 - 2  z^{- 1} + 0.75  z^{- 2}} = \frac{3  z - 3.5}{z^{2} - 2  z + 0.75} = \frac{3  z - 3.5}{(z - 0.5)  (z - 1.5)}
}{= \frac{2}{z - 0.5} + \frac{1}{z - 1.5} = \frac{2  z}{z - 0.5}  z^{- 1} + \frac{z}{z - 1.5}  z^{- 1}} \label{eq:6.94}
\end{align}
Durch Rücktransformation errechnet sich die Impulsantwort zu
\begin{align}
g[ k] &= g_{1}[ k] + g_{2}[ k] = 2  0.5^{k - 1}  \sigma[ k - 1] + 1.5^{k - 1}  \sigma[ k - 1] \label{eq:6.95}
\end{align}
In Bild \ref{fig:6.10} sind die Lage der Pole und Nullstelle sowie die Impulsantwort des Systems dargestellt.

\manualFigure
  [alt={Lage der Pol- und Nullstellen sowie und Impulsantwort bei einfachen reellen Polen}]
  {./media/image123.png}
  {Lage der Pol- und Nullstellen sowie und Impulsantwort bei einfachen reellen Polen}
  {\label{fig:6.10}}


Der Pol an der Stelle \(α_{1} = 0.5\) liegt innerhalb des Einheitskreises, der korrespondierende Teil der Impulsantwort \(g_{1}\){[}k{]} ist konvergent. Der Pol an der Stelle \(α_{2} = 1.5\) liegt außerhalb des Einheitskreises und führt zu einem divergenten Anteil \(g_{2}\){[}k{]} an der Impulsantwort. Die Summe der Impulsantworten ist damit divergent. Sobald ein Pol der Übertragungsfunktion außerhalb des Einheitskreises liegt, divergiert die Impulsantwort, und das System ist instabil. Ausführlich wird die Stabilität von Systemen in Abschnitt \ref{sec:6.3.11} diskutiert.

\end{example}

\subsection{Übertragungsfunktionen mit mehrfachen Polen}\label{sec:uxfcbertragungsfunktionen-mit-mehrfachen-polen}

Liegt ein N-facher Pol an der Stelle α ≠ 0 vor, ergibt sich für die Partialbruchzerlegung dieses Teils der Übertragungsfunktion der Ansatz
\begin{align}
G(z) &= \frac{B(z)}{(z - \alpha)^{N}} = \sum_{n = 1}^{N}\frac{A_{n}}{(z - \alpha)^{n}} \label{eq:6.96}
\end{align}
Die Rücktransformation der einzelnen Summanden kann mit Korrespondenz 9 angegeben werden zu
\begin{align}
g_{n}[ k] &= A_{n}  \left( \begin{array}{r}
k - 1 \\
n - 1
\end{array} \right)  \alpha^{k - n}  \sigma[ k - n] \label{eq:6.97}
\end{align}
Zum Beispiel ergibt sich für einen zweifachen Pol
\begin{align}
g[ k] &= \sum_{n = 1}^{2}{A_{n}  \left( \begin{array}{r}
k - 1 \\
n - 1
\end{array} \right)  \alpha^{k - n}  \sigma[ k - n] = A_{1}  \alpha^{k - 1}  \sigma[ k - 1] + A_{2}  (k - 1)  \alpha^{k - 2}  \sigma[ k - 2]} \label{eq:6.98}
\end{align}
Die Impulsantwort ist unendlich lang, es handelt sich also um ein IIR-System. Der Ausdruck
\begin{align}
\left( \begin{aligned}
 & k - 1 \\
 & n - 1
\end{aligned} \right) = \frac{(k - 1)!}{(n - 1)!  (k - 1 - n + 1)!} = \frac{(k - 1)!}{(n - 1)!  (k - n)!} \label{eq:6.99}
\end{align}
führt zu einem Polynom in k mit der Ordnung n - 1. Da die Exponentialfunktion für α ≠ 1~stärker steigt oder fällt als jede Potenz von k, ist die Stabilitätsbetrachtung von diesem Ausdruck unabhängig. Deshalb bestimmt auch in diesem Fall der Betrag des Pols α, ob die Impulsantwort gegen den Wert null konvergiert, und entscheidet damit über die Stabilität des Systems. Der Sonderfall eines Pols mit dem Betrag 1 wird in Abschnitt \ref{sec:6.3.11} diskutiert.

\begin{example}{Übertragungsfunktion mit mehrfachen reellen Polen}
\label{beispiel-uxfcbertragungsfunktion-mit-mehrfachen-reellen-polen}

Die untenstehende Übertragungsfunktion \(G(z)\) soll interpretiert werden. Ihr Zählergrad ist kleiner als der Nennergrad, und sie hat einen doppelten Pol an der Stelle α = - 0.5. Es ergibt sich die Partialbruchzerlegung
\begin{align}
G(z) &= \frac{z + 1}{(z + 0.5)^{2}} = \frac{1}{z + 0.5} + \frac{0.5}{(z + 0.5)^{2}} = \frac{z}{z + 0.5}  z^{- 1} - \frac{- 0.5  z}{(z + 0.5)^{2}}  z^{- 1} \label{eq:6.100}
\end{align}
Zur Rücktransformation werden die Korrespondenzen 5 und 6~sowie der Verschiebungssatz verwendet.
\begin{align}
g[ k] &= g_{1}[ k] + g_{2}[ k] = ( - 0.5)^{k - 1}  \sigma[ k - 1] - (k - 1)  ( - 0.5)^{k - 1}  \sigma[ k - 2] \label{eq:6.101}
\end{align}
In Bild \ref{fig:6.11} sind die Lage der Nullstelle und Pole sowie die Impulsantwort des Systems dargestellt.

\manualFigure
  [alt={Lage der Nullstelle und Pole sowie die Impulsantwort bei mehrfachem reellem Pol}]
  {./media/image124.png}
  {Lage der Nullstelle und Pole sowie die Impulsantwort bei mehrfachem reellem Pol}
  {\label{fig:6.11}}


Da der mehrfache Pol innerhalb des Einheitskreises liegt, ist das System stabil. Aufgrund des negativen Vorzeichens sind die beiden Teile der Impulsantwort alternierend, und es ergibt sich eine schwingende Impulsantwort. Besitzt die Übertragungsfunktion einen einfachen oder mehrfachen Pol mit negativem Vorzeichen,
\begin{align}
G(z) &= \frac{z}{z - r_{0}  e^{j  \pi}} \label{eq:6.102}
\end{align}
wechselt das Vorzeichen der Impulsantwort des zeitdiskreten Systems. Die Impulsantwort schwingt mit der normierten Kreisfrequenz π.
\begin{align}
g[ k] &= r_{0}^{k}  e^{j  \pi  k}  \sigma[ k] = r_{0}^{k}  ( - 1)^{k}  \sigma[ k] \label{eq:6.103}
\end{align}
Diese Systemeigenschaft besitzt im zeitkontinuierlichen Bereich kein Äquivalent.

\end{example}

\subsection{Übertragungsfunktion mit konjugiert komplexem Polpaar}\label{sec:uxfcbertragungsfunktion-mit-konjugiert-komplexem-polpaar}

Im Fall konjugiert komplexer Polpaare weist die Übertragungsfunktion Partialbrüche der Form
\begin{align}
G(z) &= \frac{A_{1}}{z - r  e^{j  \varphi}} + \frac{A_{2}}{z - r  e^{- j  \varphi}} = \frac{A}{z - r  e^{j  \varphi}} + \frac{A^{*}}{z - r  e^{- j  \varphi}} = \frac{A  e^{j  \varphi_{A}}}{z - r  e^{j  \varphi}} + \frac{A  e^{- j  \varphi_{A}}}{z - r  e^{- j  \varphi}} \label{eq:6.104}
\end{align}
auf. Damit das System reelle Koeffizienten \(a_{n}\) und \(b_{m}\) besitzt, müssen die Koeffizienten der Partialbrüche \(A_{1}\) und \(A_{2}\) ebenfalls konjugiert komplex zueinander sein. Die Impulsantwort errechnet sich zu
\begin{align}
{g\left[k\right]} &=  {A e^{j \varphi _{A} }  \left(r e^{j \varphi } \right)^{k-1}  \sigma \left[k-1\right]+A e^{-j \varphi _{A} }  \left(r e^{-j \varphi } \right)^{k-1}  \sigma \left[k-1\right]} \\
{} &=  {A r^{k-1} \left(e^{j \varphi _{A} }  e^{j \varphi  \left(k-1\right)} +e^{-j \varphi _{A} }  e^{-j \varphi  \left(k-1\right)} \right) \sigma \left[k-1\right]} \\
{} &=  {2\,\text{A} r^{k-1}  \cos \left(\varphi  \left(k-1\right)+\varphi _{A} \right) \sigma \left[k-1\right]}
 \label{eq:6.105}
\end{align}
Die Folge ist eine gedämpfte harmonische Folge. Die Phasenlage ϕ des Polpaares definiert, mit welcher Frequenz die Folge schwingt. Der Betrag der Pole r ist für das Abfallen beziehungsweise Ansteigen der Amplitude verantwortlich. Die Koeffizienten A und A* bestimmen über die Amplitude und Phasenlage der Impulsantwort. Die Impulsantwort ist wieder unendlich lang, es handelt sich also erwartungsgemäß um ein IIR-System.

\begin{example}{Übertragungsfunktion mit konjugiert komplexem Polpaar}
\label{beispiel-uxfcbertragungsfunktion-mit-konjugiert-komplexem-polpaar}

Die z-Transformierte \(G(z)\) soll in den Zeitbereich zurücktransformiert werden. Ihr Zählergrad ist kleiner als der Nennergrad, und sie hat ein konjugiert komplexes Polpaar an der Stelle α = + 0.5 ± 0.5j. Es ergibt sich die Partialbruchzerlegung
\begin{align}
G(z) &= \frac{2  z}{z^{2} - z + 0.5} = \frac{1 - j}{z - 0.5 - 0.5  j} + \frac{1 + j}{z - 0.5 + 0.5  j} = \frac{\sqrt{2}  e^{- j  \pi/4}}{z - \frac{e^{j  \pi/4}}{\sqrt{2}}} + \frac{\sqrt{2}  e^{j  \pi/4}}{z - \frac{e^{- j  \pi/4}}{\sqrt{2}}} \label{eq:6.106}
\end{align}
Die Rücktransformation führt nach den Ausführungen oben zu der Impulsantwort
\begin{align}
g[ k] &= 2  \sqrt{2}  \left( \frac{1}{\sqrt{2}} \right)^{k - 1}  \cos\left( \frac{\pi}{4}  (k - 1) - \frac{\pi}{4} \right)  \sigma[ k - 1] \label{eq:6.107}
\end{align}
In Bild \ref{fig:6.12} sind die Lage der Pole und Nullstelle sowie die Impulsantwort des Systems dargestellt.

\manualFigure
  [alt={Lage der Nullstelle und Pole sowie die Impulsantwort bei konjugiert komplexem Polpaar}]
  {./media/image126.png}
  {Lage der Nullstelle und Pole sowie die Impulsantwort bei konjugiert komplexem Polpaar}
  {\label{fig:6.12}}


Die Impulsantwort eines zeitdiskreten Systems mit konjugiert komplexem Polpaar schwingt. Liegen die Pole innerhalb des Einheitskreises nimmt die Amplitude der Schwingung ab.

\end{example}

\subsection{Übertragungsfunktion mit N-fachem Pol im Koordinatenursprung}\label{sec:uxfcbertragungsfunktion-mit-n-fachem-pol-im-koordinatenursprung}

Wenn die Übertragungsfunktion \(G(z)\) nur Pole an der Stelle \(α = 0\) besitzt, sind alle Koeffizienten\\
\(a_{n} = 0\) bis auf den Koeffizienten \(a_{N}\). Damit gilt
\begin{align}
G(z) &= \frac{\sum_{m = 0}^{M}{b_{m}  z^{m}}}{a_{N}  z^{N}} = \sum_{m = 0}^{M}{\frac{b_{m}}{a_{N}}  z^{m - N}} = \frac{b_{0}}{a_{N}}  z^{- N} + \frac{b_{1}}{a_{N}}  z^{1 - N} + ... + \frac{b_{M}}{a_{N}}  z^{M - N} \label{eq:6.108}
\end{align}
Bei kausalen Systemen ist der Zählergrad M kleiner oder gleich dem Nennergrad N. Damit sind alle Exponenten der komplexen Variablen z kleiner oder gleich null, und es kann der Verschiebungssatz der z-Transformation angewendet werden. Die Impulsantwort ergibt sich damit zu
\begin{align}
g[ k] &= \frac{b_{0}}{a_{N}}  \delta[ k - N] + \frac{b_{1}}{a_{N}}  \delta[ k - N + 1] + ... + \frac{b_{M}}{a_{N}}  \delta[ k - N + M] \label{eq:6.109}
\end{align}
Die Impulsantwort besteht aus N + 1 Impulsen. Die Koeffizienten der Impulse entsprechen den Koeffizienten der Übertragungsfunktion \(G(z)\). Im Gegensatz zu den bisher behandelten Impulsantworten ist sie nur an N + 1 Stellen von null verschieden. Deshalb wird die Impulsantwort als Finite-Impulse-Response und die Systeme als Finite-Impulse-Response-Systeme (FIR-Systeme) bezeichnet.

\begin{example}{N-facher Pol im Koordinatenursprung}
\label{beispiel-n-facher-pol-im-koordinatenursprung}

Die Übertragungsfunktion \(G(z)\) soll in den Zeitbereich zurücktransformiert werden. Sie hat einen\\
4-fachen Pol an der Stelle \(α = 0\). Das System besitzt eine endliche Impulsantwort.
\begin{align}
G(z) &= \frac{z^{4} + 2  z^{3} + 3  z^{2} + 2  z + 1}{9  z^{4}} = \frac{1}{9}  \left( 1 + 2  z^{- 1} + 3  z^{- 2} + 2  z^{- 3} + 1  z^{- 4} \right) \label{eq:6.110}
\end{align}
Mit dem Verschiebungssatz der z-Transformation ergibt sich im Zeitbereich die Folge
\begin{align}
g[ k] &= \frac{1}{9}  \delta[ k] + \frac{2}{9}  \delta[ k - 1] + \frac{3}{9}  \delta[ k - 2] + \frac{2}{9}  \delta[ k - 3] + \frac{1}{9}  \delta[ k - 4] \label{eq:6.111}
\end{align}
Die Werte der Impulsantwort sind identisch zu den Koeffizienten des Filters. In Bild \ref{fig:6.13} sind die Lage der Pole und Nullstellen sowie die Impulsantwort des Systems dargestellt.

\manualFigure
  [alt={Pollage und Impulsantwort bei Polen im Koordinatenursprung}]
  {./media/image127.png}
  {Pollage und Impulsantwort bei Polen im Koordinatenursprung}
  {\label{fig:6.13}}


Die Impulsantwort hat N + 1 = 5 von null verschiedene Werte. Für k ≥ 5 bleibt die Impulsantwort null.

\end{example}

\subsection{Übertragungsfunktion invertierbarer Systeme}\label{sec:6.3.9}

Bei der Übertragung von Signalen ist es unter Umständen erforderlich, Verzerrungen einer Übertragungsstrecke zu kompensieren. Dazu wird ein System verwendet, das ein inverses Systemverhalten aufweist. Angenommen ein Signal wird über ein System \(G_{1}(z)\) übertragen, das die Form
\begin{align}
G_{1}(z) &= \frac{z - \beta}{z - \alpha} \label{eq:6.112}
\end{align}
aufweist. Zur Kompensation wird ein System verwendet, das im Idealfall folgende Bedingung erfüllt
\begin{align}
G_{1}(z)  G_{2}(z) &= 1 \label{eq:6.113}
\end{align}
In diesem Fall würden die Verzerrungen, die durch die Signalübertragung im System \(G_{1}(z)\) entstanden sind, ideal kompensiert. Auflösen der Bedingung führt zu der Übertragungsfunktion
\begin{align}
G_{2}(z) &= \frac{1}{G_{1}(z)} = \frac{1}{\frac{z - \beta}{z - \alpha}} = \frac{z - \alpha}{z - \beta} \label{eq:6.114}
\end{align}
Die Nullstellen der Übertragungsfunktion \(G_{1}(z)\) werden zu Polen der inversen Übertragungsfunktion \(G_{2}(z)\). Das inverse System \(G_{2}(z)\) muss stabil sein, damit es realisierbar ist. Damit darf ein System, dessen Verhalten kompensiert werden soll, nur Nullstellen innerhalb des Einheitskreises besitzen. Bild \ref{fig:6.14} zeigt ein Beispiel für die Pol- und Nullstellenlage invertierbarer und nicht invertierbarer Systeme.

\begin{figure}
  \centering
  \includegraphics[alt={Pollage für ein invertierbares und ein nicht invertierbares System}]{./media/6_4_pollage_fr_ein_invertierbares_un.png}
  \caption{Pollage für ein invertierbares und ein nicht invertierbares System}
  \label{fig:6.14}
\end{figure}

Beide Systeme sind stabil, da der Pol bei beiden Systemen im Einheitskreis liegt. Das linke System ist invertierbar, da die Nullstelle innerhalb des Einheitskreises liegt. Das rechte System ist wegen der Lage der Nullstelle außerhalb des Einheitskreises nicht invertierbar.

\subsection{Zusammenfassung Interpretation der Übertragungsfunktion}\label{sec:zusammenfassung-interpretation-der-uxfcbertragungsfunktion}

Bei der Interpretation von Übertragungsfunktionen werden unterschiedliche Systemeigenschaften aufgezeigt. Tabelle \ref{tab:6.4} fasst die an der Übertragungsfunktion ablesbaren Systemeigenschaften zusammen. Der Zusammenhang zwischen Pollage und Impulsantwort ist für einfache reelle Pole und für konjugiert komplexe Polpaare in Tabelle \ref{tab:6.5} dargestellt.

\begin{summary}[label=tab:6.4]{Tabellarische Übersicht der an der Übertragungsfunktion ablesbaren Systemeigenschaften}
\begin{center}
  \begin{tabular}{ll}
    \toprule
    Eigenschaft & Übertragungsfunktion \\
    \midrule
    Kausalität & Zählergrad M ≤ Nennergrad N \\
    Sprungfähigkeit & Zählergrad M = Nennergrad N \\
    Schwingungsneigung & negative reelle Pole oder konjugiert komplexe Polpaare \\
    System mit endlicher Impulsantwort 
Finite-Impulse-Response (FIR-System)
 & Alle Pole im Koordinatenursprung \\
    System mit unendlicher Impulsantwort 
Infinite-Impulse-Response (IIR-System)
 & 
Mindestens ein Pol nicht 
im Koordinatenursprung \\
    Stabile invertierbare Systeme & 
Pole und Nullstellen innerhalb des Einheitskreises, 
Zählergrad M = Nennergrad N \\
    Verstärkung & Übertragungsfunktion G(\(z = 1\)) \\
    \bottomrule
  \end{tabular}
\end{center}
\end{summary}

\begin{table}
  \centering
  \caption{Zusammenhang zwischen Pollage und Impulsantwort}
  \label{tab:6.5}
  \begin{tabular}{ll}
    \toprule
    Pollage & Impulsantwort \\
    \midrule
    \includegraphics[width=1.9685in,height=1.31102in,alt={Ein Bild, das Reihe, Kreis, Screenshot, Diagramm enthält. Automatisch generierte Beschreibung}]{./media/image90.png} & \includegraphics[width=1.9685in,height=1.31102in,alt={Ein Bild, das Screenshot, Reihe enthält. Automatisch generierte Beschreibung}]{./media/image129.png} \\
    \includegraphics[width=1.9685in,height=1.31102in,alt={Ein Bild, das Reihe, Kreis, Screenshot, Diagramm enthält. Automatisch generierte Beschreibung}]{./media/image92.png} & \includegraphics[width=1.9685in,height=1.31102in,alt={Ein Bild, das Screenshot, Reihe, Rechteck, Design enthält. Automatisch generierte Beschreibung}]{./media/image130.png} \\
    \includegraphics[width=1.9685in,height=1.31102in,alt={Ein Bild, das Reihe, Screenshot, Diagramm, Kreis enthält. Automatisch generierte Beschreibung}]{./media/image93.png} & \includegraphics[width=1.9685in,height=1.31102in,alt={Ein Bild, das Reihe, Screenshot enthält. Automatisch generierte Beschreibung}]{./media/image131.png} \\
    \includegraphics[width=1.9685in,height=1.31102in,alt={Ein Bild, das Reihe, Screenshot, Diagramm, Kreis enthält. Automatisch generierte Beschreibung}]{./media/image101.png} & \includegraphics[width=1.9685in,height=1.31102in,alt={Ein Bild, das Screenshot, Reihe, Rechteck, Diagramm enthält. Automatisch generierte Beschreibung}]{./media/image132.png} \\
    \includegraphics[width=1.9685in,height=1.31102in,alt={Ein Bild, das Reihe, Diagramm, Kreis, Screenshot enthält. Automatisch generierte Beschreibung}]{./media/image103.png} & \includegraphics[width=1.9685in,height=1.31102in,alt={Ein Bild, das Reihe, Diagramm, Screenshot, Rechteck enthält. Automatisch generierte Beschreibung}]{./media/image133.png} \\
    \includegraphics[width=1.9685in,height=1.31102in,alt={Ein Bild, das Reihe, Diagramm, Screenshot, Kreis enthält. Automatisch generierte Beschreibung}]{./media/image105.png} & \includegraphics[width=1.9685in,height=1.31102in,alt={Ein Bild, das Reihe, Screenshot, Diagramm, Rechteck enthält. Automatisch generierte Beschreibung}]{./media/image134.png} \\
    \bottomrule
  \end{tabular}
\end{table}

\subsection{Stabilität und Pole der Übertragungsfunktion}\label{sec:6.3.11}

Bei der Einführung des Begriffes der Stabilität in Abschnitt \ref{sec:4.4.4} wird gezeigt, dass asymptotisch stabile Systeme eine abklingende Impulsantwort aufweisen müssen. Die Übertragungsfunktion \(G(z)\) ist die z-Transformierte der Impulsantwort \(g[k]\). Deshalb kann die Stabilitätsbetrachtung auch im\\
z-Bereich erfolgen. Liegen ausschließlich Pole mit einem Betrag \textbar \(α_{n}\)\textbar{} \textless{} 1 vor, konvergiert die Impulsantwort \(g[k]\) für k → ∞ gegen null. Das System ist asymptotisch stabil. Liegt mindestens ein Pol mit einem Betrag \textbar \(α_{n}\)\textbar{} \textgreater{} 1 vor, divergiert die Impulsantwort. Das System ist instabil. Bei einem Pol mit einem Betrag \textbar α\textbar{} = 1 entscheidet die Mehrfachheit über die Konvergenz. Zu der Übertragungsfunktion
\begin{align}
G(z) &= \frac{B(z)}{\left( z - e^{j  \varphi} \right)^{N}} = \sum_{n = 1}^{N}\frac{A_{n}}{\left( z - e^{j  \varphi} \right)^{n}} \label{eq:6.115}
\end{align}
gehört die Impulsantwort
\begin{align}
g[ k] &= \sum_{n = 1}^{N}{A_{n}  \left( \begin{array}{r}
k - 1 \\
n - 1
\end{array} \right)  e^{j  \varphi  (k - n)}  \sigma[ k - n]} \label{eq:6.116}
\end{align}
Der Binomialkoeffizient zum Beispiel für \(n = 2\) zu einem Faktor
\begin{align}
\left( \begin{array}{r}
k - 1 \\
2 - 1
\end{array} \right) &= \frac{(k - 1)!}{(2 - 1)!  (k - 2)!} = (k - 1) \label{eq:6.117}
\end{align}
Der Faktor steigt linear mit dem Folgenindex k. Er wird mit einer Exponentialfunktion multipliziert, deren Betrag konstant ist. Damit divergiert die Impulsantwort für eine Vielfachheit N ≥ 2.

\begin{example}{Stabilitätsnachweis über Pollage}
\label{beispiel-stabilituxe4tsnachweis-uxfcber-pollage}

Ein System mit der Übertragungsfunktion
\begin{align}
G(z) &= \frac{1}{z^{3} + 0.1  z^{2} - 0.25  z - 0.025} \label{eq:6.118}
\end{align}
soll auf Stabilität geprüft werden. Die Übertragungsfunktion des Systems kann in Linearfaktorschreibweise dargestellt werden.
\begin{align}
G(z) &= \frac{1}{(z + 0.5)  (z + 0.1)  (z - 0.5)} \label{eq:6.119}
\end{align}
Alle Pole weisen einen Betrag r \textless{} 1 auf, das System ist damit stabil.

\end{example}

\begin{example}{Stabilitätsprüfung eines Systems mit doppeltem Pol}
\label{beispiel-stabilituxe4tspruxfcfung-eines-systems-mit-doppeltem-pol}

Ein System mit der Übertragungsfunktion
\begin{align}
G(z) &= \frac{1}{(z - 0.2)  (z + 1)^{2}} \label{eq:6.120}
\end{align}
weist einen Pol an der Stelle \(α_{1} = 0.2\) und einen doppelten Pol an der Stelle \(α_{2}\) = - 1 auf. Es ergeben sich die Partialbruchdarstellung
\begin{align}
G(z) &= \frac{1}{(z - 0.2)  (z + 1)^{2}} = \frac{0.6944}{z - 0.2} - \frac{0.6944}{z + 1} + \frac{1.25}{(z + 1)^{2}} \label{eq:6.121}
\end{align}
und die Impulsantwort
\begin{align}
g[ k] &= 0.6944  0.2^{k - 1}  \sigma[ k - 1] - 0.6944  ( - 1)^{k - 1}  \sigma[ k - 1] + 1.25  (k - 1)  ( - 1)^{k - 1}  \sigma[ k - 1] \label{eq:6.122}
\end{align}
Wegen des letzten Summenden divergiert die Impulsantwort. Er ergibt sich aus dem doppelten Pol \(α_{2}\) mit einem Betrag von \textbar \(α_{2}\)\textbar{} = 1. Das System ist instabil.

\end{example}

Die Analyse der Pollage führt zu einer Stabilitätsbewertung, die in Tabelle \ref{tab:6.6} zusammengefasst ist. Unter Berücksichtigung der in Bild \ref{fig:5.2} gezeigten Abbildung vom Laplace-Bereich in den z-Bereich entsprechen die Stabilitätskriterien denen zeitkontinuierlicher Systeme im Laplace-Bereich.

\begin{table}
  \centering
  \caption{Zusammenhang zwischen Polen der Übertragungsfunktion und der Stabilität von zeitdiskreten LTI-Systemen, die sich über lineare Differenzengleichungen mit konstanten Koeffizienten beschreiben lassen}
  \label{tab:6.6}
  \begin{tabular}{ll}
    \toprule
    Eigenschaft & Pole \(α_{n}\) der Übertragungsfunktion \\
    \midrule
    Asymptotisch stabiles System & 
Alle Pole \(α_{n}\) besitzen einen 
Betrag \textbar \(α_{n}\)\textbar{} \textless{} 1 \\
    Grenzstabiles System & 
Alle Pole \(α_{n}\) besitzen einen 
Betrag \textbar \(α_{n}\)\textbar{} \textless{} 1, 
zusätzlich liegt mindestens ein 
einfacher Pole mit Betrag \textbar \(α_{n}\)\textbar{} = 1 vor \\
    Instabiles System & 
Es existiert mindestens ein Pol \(α_{n}\) 
mit einem Betrag \textbar \(α_{n}\)\textbar{} \textgreater{} 1 
oder ein mehrfacher Pol mit Betrag \textbar \(α_{n}\)\textbar{} = 1 \\
    \bottomrule
  \end{tabular}
\end{table}

Diese Diskussion der Pollage entspricht der Diskussion von Lösungen der charakteristischen Gleichung in Abschnitt \ref{sec:4.3.3}. Bei der Pollage werden die Pole der Übertragungsfunktion bestimmt:
\begin{align}
\alpha^{N} + a_{1}  \alpha^{N - 1} + a_{2}  \alpha^{N - 2} + ... + a_{N} &= 0 \label{eq:6.123}
\end{align}
In Abschnitt \ref{sec:4.3.3} werden die Lösungen der charakteristischen Gleichung analysiert.
\begin{align}
\lambda^{N} + a_{1}  \lambda^{N - 1} + a_{2}  \lambda^{N - 2} + ... + a_{N} &= 0 \label{eq:6.124}
\end{align}
Unter der Annahme, dass in der Übertragungsfunktion keine gemeinsamen Pole und Nullstellen auftreten, sind beide Gleichungen identisch, sodass die Diskussion der Lösungen zu identischen Aussagen führen muss. Damit haben aber auch beide Gleichungen das Problem, dass sie nur für Systeme mit einer Ordnung \(N = 3\) analytisch gelöst werden können. Für Ordnungen N \textgreater{} 3 kann eine Lösung nur numerisch bestimmt werden.
